\section{AALinker}

We propose a multi-agent workflow for recovering  trace links between architectural doc and model components.
Given a architecture document $D$ consisting of sentences ${s_1, \ldots, s_n}$ and an architectural model containing components ${c_1, \ldots, c_m}$, our approach produces a set of trace links $L \subseteq S
  \times C$, where each link $(s_i, c_j)$ asserts that sentence $s_i$ discusses component $c_j$'s architectural responsibilities.

The workflow is organized as a staged DAG with three layers
\begin{enumerate}
\item \textbf{Project knowledge discovery} (layer~1): Three agents analyze and extract structured knowledge from architectural documents and models. 
And the extracted knowledge will be reused in downstream layers.
\item \textbf{Candidate Link recovery} (layer~2): Four complementary agents recover candidate links via different mechanisms: component mentioning match, entity extraction, component coreference resolution, and deterministic partial-reference injection. 
The strategies are largely independent: each targets a different type of implicit component reference.
\item \textbf{Architectural judge review} (layer~3): A single stage merges all candidate links, judging recovered trace links with discovered project knowledge and traceability knowledge.
\todo{benefits}
\end{enumerate}

\subsection{Project Knowledge Discovery}
%%% provides initial seed links, could be transarc or ilinkerX
Layer~1 builds a shared knowledge base from the architectural document and model.
It first call two sub-agent analyze the architectural model to obtaining component relationships and indentify ambiguous components
and document to extract component alias each.
Then merge them to obtain: list of ambiguous components, alias mapping table, component relationships. 

\subsubsection{Architectural Model Understanding}\label{sec:model-understanding}
Architectural components operates on high level, so containing ambiguous as real world concepts.
Component names vary widely in specificity.
Constructed identifiers like \texttt{TaskScheduler} unambiguously name architectural roles, 
but single common-English words like \texttt{Core} or \texttt{Adapter} appear frequently in both refering specific component and general concepts.
\todo{source}
This pose challenges to those trace link recovery. may intro false positive. 
And confuse LLM about the trace link is not guided with reasoning paths. 
Therefore, we build a knowledge base at this step to help downstream trace link recovery\emph{architectural}
or \emph{ambiguous} using a two-step hybrid:

Therefore, blindly matching mentions of generic usage would cause false positive links.
motivating by xx paper, todo, cite ref, we extract knowledge from it.
Therefore, analyze component names from the input architecture model to identify ambiguous components, whose name can be used to as a general concepts: like logic, storage.
This can serve as knowledge and provide context for to guide futher LLM trace linker when evaluating the link. 

Motivated by previously work on camelcase identifiers, acronyms, \todo{cite}
We combine these rules with a few-shot LLM prompt to mark ambiguous components.
We marked coumpound names like CamelCase, acronyms as non-ambiguous, and use LLM to find others.\todo{cite for retionale}

Additionally, we also check component names with each other, to build a hiearchy.
Because XXX. arch model are hiearchical.
In order to disambiguates trace inks when a doc refer to parent component.
shared partial words, like XX, will be used for better coreference.
TODO example needed.


The analysis additionally extracts partial words among components. 
Because partial words can bring ambiguity by confusing 
\begin{enumerate}
  \item two components:  (e.g., \texttt{BasicScheduler}~$\to$~\texttt{Scheduler}), when the document only say Scheduler
  \item a component and general concepts: (todo: example)
\end{enumerate}
Such information will be feed to Abbreviations linking agents in layer to.

(\S\ref{sec:partial}).

\subsubsection{Architectural Document Understanding}
Since architectural documents contain heavy synonym, abbraviations, \todo{cite support}
this step is to find all alternative names for architectural components. 
\paragraph{Alternative name discovery.}
An LLM extracts three categories of alternative names:
\begin{itemize}
    \item \emph{Abbreviations}: short forms introduced via parenthetical definitions (e.g., ``Abstract Syntax Tree (AST)'' defines AST).
    \item \emph{Synonyms}: alternative names that unambiguously identify exactly one component. \todo{example from prompt}
    \item \emph{Partial references}: trailing words of multi-word component names used alone (e.g., ``Dispatcher'' for \texttt{EventDispatcher}).
\end{itemize}

\paragraph{Partial-reference refinement.}
Among the three alias types, partial references require the most careful handling: 
trailing-word shorthands are pervasive in architecture documents, but many trailing words are generic English words whose referential status depends on context.
Therefore, identify corect partial-reference needs the detected ambiguous component names from the model understanding agent.

After the extraction, a few-shot judge validates each proposed mapping.
Finally, this output a mapping table, which later linkers can use to XX. 


% partial references are ambiguous, Because XX
% generic names. 
% For each partial references to ambiguous terms,
%  we use only  capitalized occurences as they refer to specific usage. 
%TODO: check how partial refernces added 

Further 

\subsection{Coarse link recovery: Lissa style}
This sub-agent serves as provide initial seed links by looking at explicit mentions only.
It use two LLM pass each focus on different. 


\subsection{Context aware/ knowledge enhanced link recovery}\label{sec:recovery}
Layer~2 comprises four complementary strategies that recover candidate links.
Each targets a different type of reference---exact mentions, explicit entity mentions, pronominal references, and abbreviated names



\subsection{Trace Link knowledge based Judge}





